
近年来，全纯球面与\BookRecordIndex{T0143}格罗莫夫--威滕不变量引起了广泛兴趣，其中一个重要原因是：可以用它们
在紧\BookRecordIndex{T0358}辛流形的上同调环上定义一种新的乘法。本章将说明这一构造，并指出它与代数几何、可积
哈密顿系统和镜像对称等其他主题之间的一些联系。本书写作之时，这一领域正在许多有趣而
出人意料的方向上迅速发展，不可能逐一提及。关于\BookRecordIndex{T0229}量子上同调、\BookRecordIndex{T0019}Frobenius 代数和镜像对称的
更多资料，读者可参阅 Cox--Katz \srcite{76}、Dubrovin \srcite{94}、Manin
\srcite{286} 与 Voisin \srcite{409} 的著作。

本章首先讨论\BookRecordIndex{T0058}半正\BookRecordIndex{T0358}辛流形的（小）\BookRecordIndex{T0229}量子上同调环。作为加法群，它是取某个环 $\Lambda$ 为
系数的通常上同调，但其乘法经过了形变。在两个元素 $a,b\in H^*(M)$ 的通常杯积
$a\smile b$ 上，加进来自三点\BookRecordIndex{T0143}格罗莫夫--威滕不变量
$\operatorname{GW}_{A,3}^M(a,b,\,\cdot\,)$ 的若干“量子修正”项；需要系数环
$\Lambda$ 来区分不同同调类 $A\in H_2(M;\mathbb Z)$ 的贡献。\BookRecordIndex{T0228}量子杯积的结合律是\BookRecordIndex{T0125}分裂
公理的直接推论。\BookRecordIndex{T0228}量子杯积只使用了\BookRecordIndex{T0143}格罗莫夫--威滕不变量所包含的一小部分信息。在
第~\ref{sec:11-2} 节中，我们将说明所有这些信息都可编码在一个称为亏格零\BookRecordIndex{T0144}格罗莫夫--威滕
势的形式幂级数 $\Phi$ 中。我们还会说明，\BookRecordIndex{T0125}分裂公理如何化为 $\Phi$ 所满足的一组引人注目的
三阶二次偏微分方程，即 WDVV 方程。

第 11.3 节讨论四类重要例子：Fano \BookRecordIndex{T0169}环面流形、格拉斯曼流形、\BookRecordIndex{T0274}旗流形和 Calabi--Yau
流形。它们的\BookRecordIndex{T0229}量子上同调环是十分有趣的不变量，并分别与 \BookRecordIndex{T0050}Toda 晶格（\BookRecordIndex{T0274}旗流形情形）、
Verlinde 代数（格拉斯曼流形情形）和镜像对称猜想（Calabi--Yau 情形）有关。

接下来，我们定义从流形 $M$ 的哈密顿辛同胚群的基本群到其\BookRecordIndex{T0229}量子上同调环可逆元群的
\BookRecordIndex{T0047}Seidel 表示。作为应用，我们描述 Fano 与非 Fano \BookRecordIndex{T0169}环面流形的\BookRecordIndex{T0229}量子上同调。第 12 章在
\BookRecordIndex{T0132}弗洛尔同调的背景下还会回到这个主题。

最后一节重新考察第~\ref{sec:11-2} 节的讨论。它用\BookRecordIndex{T0144}格罗莫夫--威滕势在 $M$ 的上同调上
定义形式 \BookRecordIndex{T0020}Frobenius 流形结构；这一结构有时称为\BookRecordIndex{T0083}大量子上同调。该节还讨论 \BookRecordIndex{T0015}Dubrovin
联络，并说明它的平坦性等价于 Frobenius 乘积的结合律。

\section{小\BookRecordIndex{T0229}量子上同调环}\label{sec:11-1}

\paragraph{\BookRecordIndex{T0231}量子系数环。}
\BookRecordIndex{T0358}辛流形 $(M,\omega)$ 的小\BookRecordIndex{T0229}量子上同调环，作为加法群，是以适当环 $\Lambda$ 为系数的通常
上同调；选择这个环是为了能够定义新的环结构。两个类 $a,b\in H^*(M)$ 的\BookRecordIndex{T0228}量子杯积
$a*b$，是通常杯积 $a\smile b$ 与来自三点\BookRecordIndex{T0143}格罗莫夫--威滕不变量
$\operatorname{GW}_{A,3}^M(a,b,c)$（$c\in H^*(M)$）的若干量子修正项之和。形容词
“小”表示其定义只用到\BookRecordIndex{T0143}格罗莫夫--威滕不变量中很少的一部分信息。“大”\BookRecordIndex{T0229}量子上同调的结构
则以第~\ref{sec:11-2} 节所讨论的\BookRecordIndex{T0144}格罗莫夫--威滕势为形式，纳入所有亏格零\BookRecordIndex{T0143}格罗莫夫--威滕
不变量；依所采用的精确定义，它还可能包含所有高亏格不变量的信息（见 Manin
\srcite{286} 和 Usher \srcite{399, \S7}）。

对任意给定乘积 $a*b$，可能有无穷多个不同的同调类 $A$ 对它作出贡献。因此我们必须采用
允许上同调类按 $A\in H_2(M)$ 标号作无穷求和的系数环。选择很多，例如一个或多个变量的
多项式环、Laurent 级数环，或者由 $\omega$ 的上同调类确定的 Novikov 环。环的具体选择
既取决于底\BookRecordIndex{T0072}层\BookRecordIndex{T0358}辛流形，也取决于想在\BookRecordIndex{T0229}量子上同调的环结构中保留多少\BookRecordIndex{T0143}格罗莫夫--威滕不变量的
信息。若对某个实数 $\lambda$ 有 $c_1(TM)=\lambda[\omega]$，则 $\lambda>0$（单调
情形）、$\lambda=0$（Calabi--Yau 情形）与 $\lambda<0$ 三种情形颇为不同。粗略地说，
它们对应于代数几何中 Fano（正曲率）、椭圆（零曲率）和一般型（负曲率）之间的区别。

依循 Cieliebak--Salamon \srcite{69}，我们引入一个公理化框架，以统一处理这些不同选择。
\BookRecordIndex{T0358}辛流形 $(M,\omega)$ 的（辛）\emph{\BookRecordIndex{T0390}有效锥}定义为
\[
  K^{\mathrm{eff}}(M,\omega)
  :=\left\{A\in H_2(M)\ \middle|\
  \exists\,A_1,\ldots,A_N\in H_2(M):
  A=\sum_iA_i,\ 
  \operatorname{GW}_{A_i,3}^M\ne0\right\}.
\]
这是 $H_2(M):=H_2(M;\mathbb Z)/\text{挠子群}$ 中的一个锥，包含在 $\pi_2(M)$ 经
Hurewicz 同态的像内。因为 $\operatorname{GW}_{0,3}^M\ne0$，它包含零元。对每一个
$\omega$-\BookRecordIndex{T0372}驯服\BookRecordIndex{T0088}殆复结构 $J$，$K^{\mathrm{eff}}(M,\omega)$ 中的每个同调类都能由有限条
$J$-全纯球面表示。因此，$\omega$ 在这个锥的每个非零元素上取正值。

锥 $K^{\mathrm{eff}}(M,\omega)$ 是一个有用的概念工具，但可能很难计算。从辛几何观点
转到复几何观点时，会遇到\BookRecordIndex{T0390}有效锥的另一个自然定义；它在代数几何中起重要作用，而且有时
更易计算。

\begin{remark}\label{rem:11-1-1}
\begin{enumerate}[label=\textnormal{(\roman*)}]
\item 殆复流形 $(M,J)$ 的\emph{\BookRecordIndex{T0390}有效锥}定义为
\[
  K^{\mathrm{eff}}(M,J)
  :=\{A\in H_2(M)\mid \text{存在类 $A$ 中的 \BookRecordIndex{T0003}$J$-全纯曲线}\}.
\]
这是代数几何所用的定义。这里允许全纯曲线不连通、具有任意亏格，也包括常值曲线。若
$J$ 被\BookRecordIndex{T0363}辛形式 $\omega$ \BookRecordIndex{T0372}驯服，则
$K^{\mathrm{eff}}(M,\omega)\subset K^{\mathrm{eff}}(M,J)$，但这一包含通常不是等号
（即便把 $K^{\mathrm{eff}}(M,J)$ 换成由 $J$-全纯球面所表示的类的非负线性组合也一样）。
例如，K3 曲面或四维环面上的任意\BookRecordIndex{T0363}辛形式都满足 $K^{\mathrm{eff}}(M,\omega)=\{0\}$。
K3 曲面上的通有复结构没有全纯球面，但某些复结构有。经典研究还表明，大多数复四维环面
$M_\Lambda=\mathbb C^2/\Lambda$ 没有任何亏格的非常值紧全纯曲线；而与 Siegel 上半空间
相联系的那些环面则可嵌入射影空间，并有非平凡\BookRecordIndex{T0390}有效锥 $K^{\mathrm{eff}}(M_\Lambda,i)$，
当然其中只会出现亏格至少为一的全纯曲线。（关于辛四维环面与复四维环面的近期研究和更多
参考文献，见 Latschev--McDuff--Schlenk \srcite{231}。）

\item 设 $(M,J)$ 是无边紧复流形，复维数为 $n$。若其反典范线丛，即切丛的最高次复外幂
\[
  L:=\Lambda_{\mathbb C}^nTM=\Lambda^{n,0}TM,
\]
是 ample 的，则称 $(M,J)$ 为 \emph{Fano 流形}\srcite{170,233}。这意味着存在整数
$k>0$，使得 $L^{\otimes k}$ 的全纯截面定义了 $M$ 到射影空间的嵌入（即
$L^{\otimes k}$ 是 very ample 的）。\sourcefootnote{原书此处写作
“define an immersion”；但 very ample 的定义唯一给出到射影空间的嵌入，
故将 immersion 补正为 embedding。}此时，射影空间上 Fubini--Study 形式的拉回在 $M$
上定义一个 凯勒 形式 $\omega$，其上同调类满足 $[\omega]=kc_1(TM)$。
Nakai--Moishezon 判据断言，在当前情形下，闭复射影流形上的全纯线丛 $L\to M$ 是 ample
的，当且仅当对每个 $i$ 以及每个能由 $M$ 的复子簇表示的同调类
$\alpha\in H_{2i}(M;\mathbb Z)$，都有
\[
  \langle c_1(L)^i,\alpha\rangle>0
\]
（见 \srcite{171, Appendix~A, Theorem~5.1}）。因此，复射影曲面 $(M,J)$ 是 Fano 的，
当且仅当 $c_1(TM)^2>0$，并且 $c_1(TM)$ 在
$K^{\mathrm{eff}}(M,J)$ 的所有非零元素上取正值（见
\srcite{171, Chapter~V, Theorem~1.10}）。

\item 若闭复射影流形 $(M,J)$ 上的全纯线丛 $L\to M$ 的 $c_1(L)$ 在
$K^{\mathrm{eff}}(M,J)$ 上非负，则称 $L$ \emph{\BookRecordIndex{T0302}数值有效（NEF）}；这个名称来自代数几何
中的 numerically effective。Kleiman 的一个定理断言，这蕴含：对每个 $i$ 以及每个能由
$M$ 的复子簇表示的同调类 $\alpha\in H_{2i}(M;\mathbb Z)$，都有
$\langle c_1(L)^i,\alpha\rangle\ge0$（见 \srcite{233, Theorem~1.4.9}）。若复射影流形
$(M,J)$ 的反\BookRecordIndex{T0103}典范丛是 NEF，即 $c_1(TM)$ 在 $K^{\mathrm{eff}}(M,J)$ 上非负，则称
$(M,J)$ 为 NEF。

\item 若闭 \BookRecordIndex{T0207}凯勒 流形 $(M,\omega,J)$ 的底\BookRecordIndex{T0072}层复流形是 Fano 的，则称它为 Fano 流形。
此时可用一条 凯勒 形式路径把 $\omega$ 连到一个属于 $c_1(TM)$ 的形式。由
注~\ref{rem:7-3-9}，\BookRecordIndex{T0143}格罗莫夫--威滕不变量在这种形变下不变。因此在许多用途下，可以假定
Fano \BookRecordIndex{T0207}凯勒 流形是单调的。（Fano \BookRecordIndex{T0169}环面流形见第 11.3.1 小节。）

\item 把这些条件推广到辛几何背景有几种合理方式。可以把辛 Fano 流形 $(M,\omega)$ 定义
为这样一种流形：$\omega$ 能经一条\BookRecordIndex{T0363}辛形式路径连到属于 $c_1(TM)$ 的\BookRecordIndex{T0363}辛形式 $\omega'$。
也可以要求 $c_1(TM)$ 在锥 $K^{\mathrm{eff}}(M,\omega)$ 或 $K^{\mathrm{eff}}(M,J)$
之一的非零球面元素上取正值。目前尚不清楚这些不同条件之间的关系。NEF 条件也有类似的
推广。特别地，要求 $c_1(TM)$ 在某个（或每个）$\omega$-\BookRecordIndex{T0372}驯服\BookRecordIndex{T0088}殆复结构 $J$ 所对应的
$K^{\mathrm{eff}}(M,J)$ 的球面元素上非负，等价于说 $J$ 在定义~\ref{def:6-4-5} 的意义下
\BookRecordIndex{T0058}半正。对这样的流形，可以用本书现有的方法定义\BookRecordIndex{T0143}格罗莫夫--威滕不变量（见
注~\ref{rem:7-3-9}）。

\item 对某些应用，可能适合扩大锥 $K^{\mathrm{eff}}(M,\omega)$，允许出现满足
$\operatorname{GW}_{A_i,k}^M\ne0$（对某个 $k$）的类 $A_i$，或允许类 $A_i$ 的某个
高亏格不变量非零。另一种可能是把 $K^{\mathrm{eff}}$ 定义为 $\pi_2(M)$ 的子集，并把
$\operatorname{GW}_A^M$ 解释为固定同伦类中 \BookRecordIndex{T0003}$J$-全纯曲线的\BookRecordIndex{T0143}格罗莫夫--威滕不变量。
不过，除非特别说明，我们都采用上面给出的定义。
\end{enumerate}
\end{remark}

下一个引理概括\BookRecordIndex{T0390}有效锥的基本性质。证明是一道简单习题；它使用重要的有限性结果
推论~\ref{cor:5-3-2}。

\begin{lemma}\label{lem:11-1-2}
设 $(M,\omega)$ 是闭\BookRecordIndex{T0358}辛流形。则\BookRecordIndex{T0390}有效锥
$K^{\mathrm{eff}}:=K^{\mathrm{eff}}(M,\omega)$ 是含零元且包含于 Hurewicz 同态像内的
加法半群。此外，对每个 $A\in K^{\mathrm{eff}}$ 都有 $\omega(A)\ge0$，且等号成立当且
仅当 $A=0$。对每个 $A\in H_2(M)$，集合
\[
  \{B\in K^{\mathrm{eff}}\mid A-B\in K^{\mathrm{eff}}\}
\]
是有限集。
\end{lemma}

设 $R$ 是有单位元的交换环。实际中，$R$ 通常取复数的标准子环
$\mathbb Z,\mathbb Q,\mathbb R,\mathbb C$ 之一；应把它看作基本（或底）系数环。也可以模
$2$ 工作，不过这在弗洛尔上同调而非\BookRecordIndex{T0229}量子上同调中更常见。考虑 $2\mathbb Z$-分次交换环
\[
  \Gamma(M,\omega):=
  \left\{\lambda:K^{\mathrm{eff}}(M,\omega)\to R\ \middle|\
  \sup_{\lambda(A)\ne0}|c_1(A)|<\infty\right\}.
\]
对 $k\in2\mathbb Z$，次数 $k$ 部分 $\Gamma^k\subset\Gamma:=\Gamma(M,\omega)$ 由所有
映射 $\lambda:K^{\mathrm{eff}}\to R$ 组成，其中 $\lambda(A)$ 只可能在
$2c_1(A)=k$ 的类 $A\in K^{\mathrm{eff}}$ 上非零。$\Gamma$ 上的加法显然定义，乘法定义为
\[
  (\lambda\mu)(A):=\sum_{B\in K^{\mathrm{eff}}}\lambda(B)\mu(A-B),
  \qquad \lambda,\mu\in\Gamma.
\]
由引理~\ref{lem:11-1-2}，右端的和有限。这个乘法保持分次：若
$\lambda\in\Gamma^k$ 且 $\mu\in\Gamma^\ell$，则
$\lambda\mu\in\Gamma^{k+\ell}$。单位元是映射 $\delta:K^{\mathrm{eff}}\to R$，其中
$\delta(0):=1$，而当 $A\ne0$ 时 $\delta(A):=0$。

\begin{definition}\label{def:11-1-3}
设 $(M,\omega)$ 是紧\BookRecordIndex{T0358}辛流形，$R$ 是有单位元的交换环。$(M,\omega)$ 的一个以 $R$ 为底环
的\emph{\BookRecordIndex{T0231}量子系数环}是三元组 $(\Lambda,\phi,\iota)$，其中 $\Lambda$ 是有单位元的
$2\mathbb Z$-分次交换环，同时也是 $R$-模；
$\phi:\Gamma(M,\omega)\to\Lambda$ 是保持 $R$-模结构和分次的环同态；
$\iota:\Lambda\to R$ 是一个 $R$-模同态，它在非零次数元素上消失，并满足
$\iota(\phi(\lambda))=\lambda(0)$。
\end{definition}

给定定义~\ref{def:11-1-3} 中的\BookRecordIndex{T0231}量子系数环 $(\Lambda,\phi,\iota)$，对
$\lambda\in\Gamma(M,\omega)$ 使用记号
\[
  \phi(\lambda)=:\sum_{A\in K^{\mathrm{eff}}}\lambda(A)e^A.
\]
注意，由于不假定 $\phi$ 为单射，$\Lambda$ 中一个元素可能有多个这样的表达。若表达唯一，
通常称之为形式和。

一个显然的例子是取 $\Lambda:=\Gamma(M,\omega)$、$\phi:=\operatorname{id}$；它适用于
每个闭\BookRecordIndex{T0358}辛流形。然而这往往不是最佳选择，例如我们可能并不确切知道
$K^{\mathrm{eff}}(M,\omega)$。下面列出其他一些更有趣的例子。注意，对每个
$A\in K^{\mathrm{eff}}$，元素 $e^A\in\Lambda$ 表示函数
$\delta_A:K^{\mathrm{eff}}\to R$ 在 $\phi$ 下的像，其中 $\delta_A(A):=1$，而当
$B\ne A$ 时 $\delta_A(B):=0$。它的次数为
$\deg(e^A)=\deg(\delta_A)=2c_1(A)$。因为 $\phi$ 是环同态，这些元素满足关系
\[
  e^{A+B}=e^Ae^B.
\]
定义~\ref{def:11-1-3} 的假设保证，只要 $\lambda\in\Gamma(M,\omega)$，和
$\lambda:=\sum_A\lambda(A)e^A$ 就是 $\Lambda$ 中良定义的元素。这些和可能是无穷和；但若
$(M,\omega)$ 单调，则它们都是有限和。

\begin{example}\label{ex:11-1-4}
\begin{enumerate}[label=\textnormal{(\roman*)}]
\item 若 $(M,\omega)$ 单调，则
\begin{equation}\label{eq:11-1-1}
  \begin{gathered}
  A\in K^{\mathrm{eff}}(M,\omega),\quad A\ne0
  \quad\Longrightarrow\quad c_1(A)>0,\\
  \#\{A\in K^{\mathrm{eff}}(M,\omega)\mid c_1(A)\le c\}<\infty
  \qquad\forall c>0.
  \end{gathered}
\end{equation}
若存在一个 $\omega$-\BookRecordIndex{T0372}驯服复结构 $J$ 使 $(M,J)$ 为 Fano 流形，上述结论也成立；特别地，
它对每个 Fano \BookRecordIndex{T0207}凯勒 流形 $(M,\omega,J)$ 成立（见注~\ref{rem:11-1-1}）。假定
\eqref{eq:11-1-1}，可以选取单变量（次数为二）多项式环
\[
  \Lambda:=R[q].
\]
同态 $\phi:\Gamma(M,\omega)\to R[q]$ 由 $e^A:=q^{c_1(A)}$ 给出；等价地，
\[
  \phi(\lambda):=\sum_{A\in K^{\mathrm{eff}}}\lambda(A)q^{c_1(A)}.
\]
同态 $\iota:R[q]\to R$ 把多项式送到 $1=q^0$ 的系数。类似的例子是
$\Lambda:=R[q,q^{-1}]$，即关于 $q,q^{-1}$ 的 Laurent 多项式环；还可以取关于 $q$ 或
$q^{-1}$ 的 Laurent 级数环。若 $R$ 是域，则这两种 Laurent 级数环也都是域。注意在
Laurent 级数的例子中，$\Lambda$ 不是各 $\Lambda^k$ 的直和。

\item 假定~\eqref{eq:11-1-1}。可以选取 $H_2(M)$ 在 $R$ 上的群环作为 $\Lambda$；于是
$\Lambda$ 的元素是如下有限形式和：
\begin{equation}\label{eq:11-1-2}
  \lambda=\sum_{A\in H_2(M)}\lambda(A)e^A,
  \qquad \lambda(A)\in R.
\end{equation}
由~\eqref{eq:11-1-1}，$\Gamma=\Gamma(M,\omega)$ 是 $\Lambda$ 的子环，且
$\phi:\Gamma\to\Lambda$ 是自然包含。映射 $\iota:\Lambda\to R$ 由
$\iota(\lambda):=\lambda(0)$ 给出。

$H_2(M)$ 的一个整数基 $A_1,\ldots,A_m$ 给出其群环与多项式环
\[
  \Lambda:=R[q_1,\ldots,q_m,q_1^{-1},\ldots,q_m^{-1}]
\]
之间的同构，其中变量 $q_i,q_i^{-1}$ 满足
\[
  \deg(q_i)=2c_1(A_i).
\]
在此记号下，同态 $\Gamma\to\Lambda$ 由 $e^{A_i}:=q_i$ 给出，同态
$\iota:\Lambda\to R$ 则取常数项。若再假定 $K^{\mathrm{eff}}(M,\omega)$ 的每个元素都能
写成 $A=\sum_i d_iA_i$，其中系数 $d_i$ 非负，那么也可以改取
\[
  \Lambda:=R[q_1,\ldots,q_m],
\]
而 $\phi$ 与 $\iota$ 仍如上定义。也可以通过减少或增加变量 $q_i$ 的数目来调整这个例子。
令 $q_i:=q^{c_1(A_i)}$，便重新得到 (i) 的例子。

\item 现在设 $(M,\omega)$ 是任意闭\BookRecordIndex{T0358}辛流形，并定义
\begin{equation}\label{eq:11-1-3}
  \Lambda_\omega:=
  \left\{\sum_{A\in H_2(M)}\lambda(A)e^A\ \middle|\
  \#\left\{A\ \middle|\ \lambda(A)\ne0,\ \omega(A)\le c\right\}<\infty
  \ \forall c>0\right\}.
\end{equation}
这是 $H_2(M)$ 的群环的一个完备化，称为 $\omega$ 的 \emph{Novikov 环}。常用定义有几个
略有不同的版本。例如，可以对球面同调群 $H_2^S(M)$ 求和；也可以对 $H_2(M)$ 中的等价类
求和，其中 $A\sim B$ 当且仅当 $\omega(A)=\omega(B)$ 且 $c_1(A)=c_1(B)$。按上面的
定义，$\Gamma(M,\omega)$ 是 $\Lambda_\omega$ 的子环，$\phi$ 是自然包含，并且
$\iota(\lambda):=\lambda(0)$。若同态
$H_2(M)\to\mathbb R:A\mapsto\omega(A)$ 是单射且 $R$ 是域，则 $\Lambda_\omega$ 也是域
（见 Hofer--Salamon \srcite{180}）。

与 (ii) 一样，$H_2(M)$ 的一个基把 $\Lambda_\omega$ 等同于由变量 $q_i,q_i^{-1}$ 中
满足有限性条件的所有形式幂级数
$\lambda=\sum_{d\in\mathbb Z^m}\lambda_dq^d$ 所组成的环：对每个 $c\in\mathbb R$，
\[
  \#\{d\in\mathbb Z^m\mid \lambda_d\ne0,\ \omega(A_d)\le c\}<\infty,
  \qquad A_d:=\sum_i d_iA_i.
\]
同态 $\phi,\iota$ 仍由 $e^{A_i}:=q_i$ 和 $\iota(\lambda):=\lambda_0$ 给出。若对所有 $i$
都有 $\omega(A_i)>0$，而 $K^{\mathrm{eff}}$ 的每个元素都是 $A_i$ 的非负线性组合，则可以
把这个例子改为只取 $q_1,\ldots,q_m$ 中的形式幂级数环。此时 $\Lambda$ 的每个元素都自动
满足有限性条件。

要强调的是，除非 $c_1(TM)$ 是挠类，$\Lambda_\omega$ 并不是 $2\mathbb Z$-分次的。否则
$\Lambda_\omega$ 不是加法子群 $\Lambda_\omega^k$（即所有满足 $2c_1(A)=k$ 的类 $A$
上的形式和）的直和，因而不是定义~\ref{def:11-1-3} 意义下的\BookRecordIndex{T0231}量子系数环。为了恢复
$2\mathbb Z$-分次，可以取 $\Lambda_\omega^k$ 的直和，即
\begin{equation}\label{eq:11-1-4}
  \Lambda_{\omega,c}:=
  \left\{\lambda\in\Lambda_\omega\ \middle|\
  \sup_{\lambda(A)\ne0}|c_1(A)|<\infty\right\}.
\end{equation}
这是其零次子环 $\Lambda_\omega^0$ 上的模，并有自然的 $2\mathbb Z$-分次；一般而言它不是
域。另一种办法是把 $\Lambda_\omega$ 看作无分次\BookRecordIndex{T0231}量子系数环。这仍然有用，但必须记住，
此时相应的\BookRecordIndex{T0229}量子上同调或\BookRecordIndex{T0132}弗洛尔同调群只能是 $\mathbb Z_2$-分次的。

\item 假定 $(M,\omega)$ 是 Calabi--Yau 流形，因而切丛的第一陈类在 $\pi_2(M)$ 上消失。
此时可取 $\Lambda=\Lambda^0:=\Lambda^{\mathrm{univ}}$ 为由如下形式幂级数组成的
\emph{\BookRecordIndex{T0270}普适 Novikov 环}：
\[
  \lambda=\sum_{\varepsilon\in\mathbb R}\lambda_\varepsilon t^\varepsilon,
  \qquad
  \#\{\varepsilon\in\mathbb R\mid
  \lambda_\varepsilon\ne0,\ \varepsilon\le c\}<\infty
  \quad\forall c\in\mathbb R.
\]
系数 $\lambda_\varepsilon$ 属于底环 $R$。若 $R$ 是域，\BookRecordIndex{T0270}普适 Novikov 环也是域。在这个例子中，
同态 $\phi$ 定义为
\begin{equation}\label{eq:11-1-5}
  \phi(\lambda):=\sum_{A\in K^{\mathrm{eff}}}\lambda(A)t^{\omega(A)},
\end{equation}
而对 $\lambda=\sum_\varepsilon\lambda_\varepsilon t^\varepsilon
\in\Lambda^{\mathrm{univ}}$，令 $\iota(\lambda):=\lambda_0$。

\item 对一般\BookRecordIndex{T0358}辛流形 $(M,\omega)$，可以结合 (i) 与 (iv)，取
$\Lambda:=\Lambda^{\mathrm{univ}}[q,q^{-1}]$ 为系数在\BookRecordIndex{T0270}普适 Novikov 环中的单变量 Laurent
多项式环，其中 $q$ 的次数为二，并定义
\[
  \phi(\lambda):=\sum_{A\in K^{\mathrm{eff}}}
  \lambda(A)t^{\omega(A)}q^{c_1(A)}.
\]
类似地，也可以在 (iii) 中取 $\Lambda=\Lambda_\omega[q,q^{-1}]$，并赋每个元素 $e^A$
以零次，从而恢复分次。

\item 若愿意采用模 $2$ 分次，则即使第一陈类在 $\pi_2(M)$ 上不消失，也可以使用\BookRecordIndex{T0270}普适
Novikov 环（或把 $\Lambda_\omega$ 看作零次环）。许多作者，例如 \srcite{128}，采用这种
做法。Usher 在 \srcite{399, \S7} 中定义了一个“普适\BookRecordIndex{T0231}量子系数环” $R_M$，它可以特化为
上面定义的许多环。

\item 例如，采用实系数并取 $\Lambda=\mathbb R$ 会很有意思。人们或许想固定某个正实数
$t$，并用~\eqref{eq:11-1-5} 定义 $\phi$。但这样一来，必须把 $\Gamma$ 限制为只包含那些
使级数~\eqref{eq:11-1-5} 绝对收敛的映射
$\lambda:K^{\mathrm{eff}}\to\mathbb R$。下面将看到，在我们的应用中，系数 $\lambda(A)$
会是形如 $\operatorname{GW}_{A,3}^M(a,b,c)$ 的某些\BookRecordIndex{T0143}格罗莫夫--威滕不变量。这就引出了一个
有趣问题：级数
\[
  \sum_{A\in H_2(M)}\operatorname{GW}_{A,3}^M(a,b,c)t^{\omega(A)}
\]
是否对充分小的 $t>0$ 收敛。这似乎是一个困难的分析问题；在本书写作之时，它仍未解决。
这个问题在 Calabi--Yau 流形和镜像对称的背景下尤其重要。（关于处理该问题的方法，见
Cox--Katz \srcite{76, \S1.3} 的进一步讨论。）
\end{enumerate}
\end{example}

\begin{remark}\label{rem:11-1-5}
记号 $e^A:=\phi(\delta_A)\in\Lambda$ 仅对 $A\in K^{\mathrm{eff}}(M,\omega)$ 有意义。不过，
在大多数例子中，同态 $\phi:\Gamma(M,\omega)\to\Lambda$ 自然延拓到 $H_2(M)$ 的群环，
此时对每个 $A\in H_2(M)$，记号 $e^A:=\phi(\delta_A)\in\Lambda$ 都有意义。（某些情形
还需限制 $c_1(A)\ge0$。）在例~\ref{ex:11-1-4}(ii)、(iii) 中，
\[
  \iota(e^A)=
  \begin{cases}1,&A=0,\\0,&\text{其他情形},\end{cases}
\]
而在例~\ref{ex:11-1-4}(i)、(iv)、(v) 中，
\[
  \iota(e^A)=
  \begin{cases}1,&c_1(A)=\omega(A)=0,\\0,&\text{其他情形}.
  \end{cases}
\]
若把注意力限制在 $A\in K^{\mathrm{eff}}(M,\omega)$，这两个公式一致。
\end{remark}

\begin{remark}\label{rem:11-1-6}
Novikov \srcite{301} 最早在闭 $1$-形式的 Morse 理论背景下引入了例~\ref{ex:11-1-4}(iii)
中 $\Lambda_\omega$ 形式的环。在这种情形，$\pi_2(M)$（在我们的表述中相应为 $H_2(M)$）
由基本群 $\pi_1(M)$ 代替，而决定有限性条件~\eqref{eq:11-1-3} 的赋值
$\pi_1(M)\to\mathbb R$ 则由一个闭 $1$-形式 $\omega$ 诱导。

正如第 12 章将看到的，Novikov 环在\BookRecordIndex{T0132}弗洛尔同调中由 $M$ 的自由环路空间上的闭 $1$-形式
（辛作用量）产生。Hofer 与 Salamon 在 \srcite{180} 中利用这一点证明了\BookRecordIndex{T0058}半正情形的
Arnold 猜想；这也显示出\BookRecordIndex{T0229}量子上同调与\BookRecordIndex{T0132}弗洛尔同调之间的紧密联系。
\end{remark}
\paragraph{\BookRecordIndex{T0229}量子上同调。}
设 $(M,\omega)$ 是闭\BookRecordIndex{T0058}半正\BookRecordIndex{T0358}辛流形，$R$ 是有单位元的交换环，而
$(\Lambda,\phi,\iota)$ 是定义~\ref{def:11-1-3} 中以 $R$ 为底环的\BookRecordIndex{T0231}量子系数环。回忆
$H^*(M)$ 表示整数上同调群 $H^*(M;\mathbb Z)$ 对挠子模的商。$(M,\omega)$ 的
\emph{以 $\Lambda$ 为系数的\BookRecordIndex{T0229}量子上同调}定义为
\[
  \operatorname{QH}^*(M;\Lambda):=H^*(M)\otimes_{\mathbb Z}\Lambda.
\]
这是带单位元的分次 $R$-模。对 $k\in\mathbb Z$，$\operatorname{QH}^*(M;\Lambda)$ 的
$k$ 次部分为
\[
  \operatorname{QH}^k(M;\Lambda)
  =\bigoplus_iH^i(M)\otimes_{\mathbb Z}\Lambda^{k-i}.
\]
$\operatorname{QH}^*(M;\Lambda)$ 的每个元素都能写成有限和
$a=\sum_i a_i\otimes\lambda_i$，其中 $a_i\in H^*(M)$、$\lambda_i\in\Lambda$。\BookRecordIndex{T0229}量子上同调
有显然的 $\Lambda$-模结构，而 $\operatorname{QH}^k(M;\Lambda)$ 是 $\Lambda^0$-模。
在某些情形，$\operatorname{QH}^*(M;\Lambda)$ 不等于所有
$k\in\mathbb Z$ 的 $\Lambda^0$-模 $\operatorname{QH}^k(M;\Lambda)$ 的直和（见
例~\ref{ex:11-1-4}(i)）；当 $\Lambda$ 具有相应性质时，它才等于该直和。

通常上同调 $H^*(M)$ 经 $a\mapsto a\otimes1$ 嵌入 $\operatorname{QH}^*(M;\Lambda)$。
还要注意，$R$-模同态 $\iota:\Lambda\to R$ 诱导一个\emph{相交配对}
\[
  \operatorname{QH}^*(M;\Lambda)\otimes
  \operatorname{QH}^*(M;\Lambda)\longrightarrow R,
\]
它在通常上同调上的限制就是标准配对（取值于 $\mathbb Z\subset R$）。具体地，对
$a=\sum_i a_i\otimes\lambda_i$ 和 $b=\sum_jb_j\otimes\mu_j$，它由
\begin{equation}\label{eq:11-1-6}
  \langle a,b\rangle
  :=\sum_{i,j}\iota(\lambda_i\mu_j)\int_Ma_i\smile b_j
\end{equation}
给出。配对~\eqref{eq:11-1-6} 未必非退化（这里非退化指：若对所有 $b$ 都有
$\langle a,b\rangle=0$，则 $a=0$）。当配对
$\Lambda\times\Lambda\to R:(\lambda,\mu)\mapsto\iota(\lambda\mu)$ 非退化时，它是
非退化的；参见习题~\ref{ex:11-1-21}。此外，配对~\eqref{eq:11-1-6} 是分次交换的：
若 $a,b\in\operatorname{QH}^*(M;\Lambda)$ 次数纯齐，则
\[
  \langle b,a\rangle=(-1)^{\deg(a)\deg(b)}\langle a,b\rangle.
\]

\begin{remark}\label{rem:11-1-7}
设 $K^{\mathrm{eff}}\to H^*(M)\otimes_{\mathbb Z}R:A\mapsto a_A$ 是任意满足
$\sup_{a_A\ne0}|c_1(A)|<\infty$ 的函数。则可能为无穷和的
\begin{equation}\label{eq:11-1-7}
  a=\sum_Aa_A\otimes e^A
\end{equation}
自然地表示 $\operatorname{QH}^*(M;\Lambda)$ 中的一个元素。为看出这一点，选取
$H^*(M)$ 的整数基 $e_0,\ldots,e_N$，写成
$a_A:=\sum_{\nu=0}^Ne_\nu\otimes\lambda_\nu(A)$，其中 $\lambda_\nu(A)\in R$。依定义，
函数 $\lambda_\nu:K^{\mathrm{eff}}\to R$ 属于环 $\Gamma(M,\omega)$，因而把作为
$A\in K^{\mathrm{eff}}$ 上无穷和给出的类 $a$ 解释为有限和
\[
  a=\sum_{A\in K^{\mathrm{eff}}}a_A\otimes e^A
  :=\sum_{\nu=0}^Ne_\nu\otimes\phi(\lambda_\nu).
\]
读者可以验证，这个定义与基的选择无关。

现在假定 $R$ 是域 $\mathbb Q,\mathbb R,\mathbb C$ 之一，并且 $\Lambda$ 是 $H_2(M)$ 的
群环或 Novikov 环 $\Lambda_\omega$（见例~\ref{ex:11-1-4}(ii)、(iii)）。此时
$\operatorname{QH}^*(M;\Lambda)$ 的每个元素都是形如~\eqref{eq:11-1-7} 的有限或可能
无穷形式和，其中 $a_A\in H^*(M;R)$，求和遍历所有 $A\in H_2(M)$。在 Novikov 情形，
该和对每个 $c\in\mathbb R$ 都满足有限性条件
\[
  \#\{A\in H_2(M)\mid a_A\ne0,\ \omega(A)\le c\}<\infty.
\]
\end{remark}

\paragraph{\BookRecordIndex{T0228}量子杯积。}
$\operatorname{QH}^*(M;\Lambda)$ 上的环结构，也称\BookRecordIndex{T0228}量子杯积，是 $\Lambda$-模的双线性
同态
\[
  \operatorname{QH}^*(M;\Lambda)\otimes
  \operatorname{QH}^*(M;\Lambda)
  \longrightarrow\operatorname{QH}^*(M;\Lambda),
  \qquad(a,b)\longmapsto a*b.
\]
由于 $\operatorname{QH}^*(M;\Lambda)$ 作为 $\Lambda$-模由 $H^*(M)$ 中的元素生成，
\BookRecordIndex{T0228}量子杯积完全由它在 $a,b\in H^*(M)$ 上的值决定。为给这些元素定义环结构，选取
$H^*(M)$（作为 $\mathbb Z$-模）的一个整数基 $e_0,\ldots,e_N$，使
$e_0=1\in H^0(M)$，且每个 $e_\nu$ 次数纯齐。定义整数矩阵
\[
  g_{\nu\mu}:=\int_Me_\nu\smile e_\mu,
\]
并用 $g^{\nu\mu}$ 表示其逆矩阵。于是，次数纯齐的两个类 $a,b\in H^*(M)$ 的\BookRecordIndex{T0228}量子杯积
定义为
\begin{equation}\label{eq:11-1-8}
  a*b:=\sum_{A\in K^{\mathrm{eff}}}\sum_{\nu,\mu}
  \operatorname{GW}_{A,3}^M(a,b,e_\nu)g^{\nu\mu}e_\mu\otimes e^A.
\end{equation}
右端与三元组 $A,\nu,\mu$ 对应的项只有在
\[
  \deg(e_\nu)+\deg(e_\mu)=\dim M,\qquad
  \dim M+2c_1(A)=\deg(a)+\deg(b)+\deg(e_\nu)
\]
时才可能非零。因此得到维数条件
\[
  \deg(a)+\deg(b)-\dim M\le2c_1(A)\le\deg(a)+\deg(b).
\]
利用注~\ref{rem:11-1-7} 可知，\eqref{eq:11-1-8} 的右端是
$\operatorname{QH}^*(M;\Lambda)$ 中良定义的元素。还要注意
$\deg(a*b)=\deg(a)+\deg(b)$。量子积也可方便地写成
\[
  a*b=\sum_{A\in K^{\mathrm{eff}}}(a*b)_A\otimes e^A,\qquad
  (a*b)_A:=\sum_{\nu,\mu}
  \operatorname{GW}_{A,3}^M(a,b,e_\nu)g^{\nu\mu}e_\mu,
\]
其中 $e^A$ 的系数
\[
  (a*b)_A\in H^{\deg(a)+\deg(b)-2c_1(A)}(M)
\]
由下列条件刻画：
\begin{equation}\label{eq:11-1-9}
  \int_M(a*b)_A\smile c
  :=\operatorname{GW}_{A,3}^M(a,b,c),
  \qquad c\in H^*(M).
\end{equation}

\begin{remark}\label{rem:11-1-8}
采用注~\ref{rem:11-1-7} 的记号，两个\BookRecordIndex{T0229}量子上同调类
$a=\sum_Aa_A\otimes e^A$ 和 $b=\sum_Bb_B\otimes e^B$ 的量子积是
\[
  a*b=\sum_{A,B,C,\nu,\mu}
  \operatorname{GW}_{C,3}^M(a_A,b_B,e_\nu)g^{\nu\mu}e_\mu
  \otimes e^{A+B+C},
\]
而配对~\eqref{eq:11-1-6} 为
\[
  \langle a,b\rangle
  =\sum_{A,B}\iota(e^{A+B})\int_Ma_A\smile b_B
  =\alpha(a*b),
  \qquad
  \alpha(a):=\sum_A\iota(e^A)\int_Ma_A.
\]
这里不要求类 $a_A$ 次数纯齐，并把 $M$ 上的积分理解为取 $2n$ 次分量后积分。（记号
$\iota(e^A)$ 见注~\ref{rem:11-1-5}。）
\end{remark}

\begin{remark}\label{rem:11-1-9}
上同调类 $(a*b)_A\in H^*(M)$ 有如下几何解释。分别设
$f_0:X_0\to M$ 和 $f_1:X_1\to M$ 是与 $a,b$ Poincaré 对偶的\BookRecordIndex{T0337}伪循环。由
\eqref{eq:11-1-9}，$a,b$ 的量子积的 $A$-分量
\[
  (a*b)_A\in H^{\deg(a)+\deg(b)-2c_1(A)}(M)
\]
可由类 $A$ 中所有经过 $f_0,f_1$ 之像的单纯 $J$-全纯球面所成的\BookRecordIndex{T0337}伪循环表示。其定义域是
满足 $u(0)=f_0(x_0)$、$u(1)=f_1(x_1)$ 的三元组
\[
  (x_0,x_1,u)\in X_0\times X_1\times\mathcal M^*(A;J)
\]
之集，\BookRecordIndex{T0337}伪循环则是\BookRecordIndex{T0283}求值映射 $(x_0,x_1,u)\mapsto u(\infty)$。对通有三元组
$(f_0,f_1,J)$，该定义域是维数
$2n+2c_1(A)-\deg(a)-\deg(b)$ 的流形，而\BookRecordIndex{T0283}求值映射是\BookRecordIndex{T0337}伪循环。由~\eqref{eq:11-1-9}，
这个\BookRecordIndex{T0337}伪循环与上同调类 $(a*b)_A$ Poincaré 对偶。

例如，考虑射影空间 $M:=\mathbb{CP}^n$，记
\[
  L:=[\mathbb{CP}^1]\in H_2(\mathbb{CP}^n),\qquad
  c:=\operatorname{PD}([\mathbb{CP}^{n-1}])\in H^2(\mathbb{CP}^n).
\]
用一个 $(n-k)$-平面 $X\subset\mathbb{CP}^n$ 表示类 $\operatorname{PD}(c^k)$，用一个
$(n-\ell)$-平面 $Y\subset\mathbb{CP}^n$ 表示 $\operatorname{PD}(c^\ell)$。假定
$k+\ell>n$ 且 $X\cap Y=\varnothing$。上同调类 $(c^k*c^\ell)_L$ Poincaré 对偶于如下
同调类：它由位于同时经过 $X,Y$ 的直线上的点所成之集表示。这个集合是复维数
$2n-k-\ell+1$ 的射影线性子空间，故
\[
  (c^k*c^\ell)_L=c^{k+\ell-n-1}.
\]
这与下面例~\ref{ex:11-1-12} 的讨论一致。
\end{remark}

\begin{remark}\label{rem:11-1-10}
两个类 $a,b\in H^*(M)$ 的杯积可写成
\[
  a\smile b=\sum_{\nu,\mu}
  \left(\int_Ma\smile b\smile e_\nu\right)g^{\nu\mu}e_\mu.
\]
因此，三重交数
\[
  c_{ijk}:=\int_Me_i\smile e_j\smile e_k
\]
可看作上同调环的结构常数。用结构常数表示时，杯积的结合律等价于
\begin{equation}\label{eq:11-1-10}
  \sum_{\nu,\mu}c_{ij\nu}g^{\nu\mu}c_{\mu k\ell}
  =(-1)^{|e_i|(|e_j|+|e_k|)}
  \sum_{\nu,\mu}c_{jk\nu}g^{\nu\mu}c_{\mu i\ell},
\end{equation}
其中用了缩写 $|e|:=\deg(e)$。同样，可以把\BookRecordIndex{T0143}格罗莫夫--威滕不变量
\[
  c_{ijk}(A):=\operatorname{GW}_{A,3}^M(e_i,e_j,e_k)
\]
看作\BookRecordIndex{T0229}量子上同调的结构常数；有一个与其结合律等价的类似方程。这解释了
第~\ref{sec:11-2} 节所讨论的 WDVV 方程为何具有那种形式。
\end{remark}

下一个命题概括量子积的主要性质。结合律最早由 Ruan--Tian \srcite{345,346} 证明，后来
Liu \srcite{248} 和本书第一版又给出了证明。它是\BookRecordIndex{T0256}粘合定理的直接推论。

\begin{proposition}\label{prop:11-1-11}
\begin{enumerate}[label=\textnormal{(\roman*)}]
\item \BookRecordIndex{T0228}量子杯积对加法满足分配律，并在下述意义下分次交换：对次数纯齐的
$a,b\in H^*(M)$，
\[
  b*a=(-1)^{\deg(a)\deg(b)}a*b.
\]
它满足结合律，并与 $\Lambda$ 的作用交换。

\item \BookRecordIndex{T0228}量子杯积的首项是标准杯积，即对所有
$a,b\in\operatorname{QH}^*(M;\Lambda)$ 都有
\[
  (a*b)_0=a\smile b.
\]
此外，只要 $a\in H^0(M)$ 或 $a\in H^1(M)$，高阶项就都消失；也就是说，对所有
$b\in H^*(M)$ 都有 $a*b=a\smile b$。因此，典范生成元
$1\in H^0(M)$ 是\BookRecordIndex{T0229}量子上同调的单位元。

\item 由~\eqref{eq:11-1-6} 定义的配对
$\operatorname{QH}^*(M;\Lambda)\otimes\operatorname{QH}^*(M;\Lambda)\to R$
与\BookRecordIndex{T0228}量子杯积\BookRecordIndex{T0348}相容，即
\[
  \langle a*b,c\rangle=\langle a,b*c\rangle
\]
对 $a,b,c\in\operatorname{QH}^*(M;\Lambda)$ 成立。
\end{enumerate}
\end{proposition}

\begin{proof}
(i) 中唯一非平凡的断言是结合律。它由第 10 章建立的\BookRecordIndex{T0143}格罗莫夫--威滕不变量\BookRecordIndex{T0125}分裂公理
（定理~\ref{thm:7-5-10}）\BookRecordIndex{T0326}推出。对四个\BookRecordIndex{T0150}固定标记点的不变量，定理~\ref{thm:7-5-10} 断言：
令 $I:=\{1,2,3,4\}$，则
\[
  \operatorname{GW}_{A,4}^{M,I}(a_1,a_2,a_3,a_4)
  =\sum_{A_0+A_1=A}\sum_{\nu,\mu}
  \operatorname{GW}_{A_0,3}^M(a_1,a_2,e_\nu)g^{\nu\mu}
  \operatorname{GW}_{A_1,3}^M(e_\mu,a_3,a_4).
\]
因此，对任意三个次数纯齐的上同调类 $a_1,a_2,a_3\in H^*(M)$ 和每个
$A\in K^{\mathrm{eff}}$，有
\[
\begin{aligned}
  ((a_1*a_2)*a_3)_A
  &=\sum_B((a_1*a_2)_B*a_3)_{A-B}\\
  &=\sum_B\sum_{\nu,\mu}
    \operatorname{GW}_{B,3}^M(a_1,a_2,e_\nu)g^{\nu\mu}
    (e_\mu*a_3)_{A-B}\\
  &=\sum_B\sum_{\nu,\mu}\sum_{i,j}
    \operatorname{GW}_{B,3}^M(a_1,a_2,e_\nu)g^{\nu\mu}
    \operatorname{GW}_{A-B,3}^M(e_\mu,a_3,e_i)g^{ij}e_j\\
  &=\sum_{i,j}\operatorname{GW}_{A,4}^{M,\{1,2,3,4\}}
    (a_1,a_2,a_3,e_i)g^{ij}e_j\\
  &=\varepsilon\sum_{i,j}\operatorname{GW}_{A,4}^{M,\{1,2,3,4\}}
    (a_2,a_3,a_1,e_i)g^{ij}e_j\\
  &=\varepsilon((a_2*a_3)*a_1)_A\\
  &=(a_1*(a_2*a_3))_A,
\end{aligned}
\]
其中使用缩写 $|a|:=\deg(a)$，并且
\[
  \varepsilon:=(-1)^{|a_1|(|a_2|+|a_3|)}.
\]
这里用到了如下事实：当 $a,b,c,d\in H^*(M)$ 次数纯齐时，不变量
$\operatorname{GW}_{A,4}^{M,\{1,2,3,4\}}(a,b,c,d)$ 在其输入的置换下分次对称。这证明了
(i)。

(ii) 的第一个断言由~\eqref{eq:11-1-9} 以及
\[
  \operatorname{GW}_{0,3}^M(a,b,c)=\int_Ma\smile b\smile c
\]
\BookRecordIndex{T0326}推出。为证明第二个断言，需要说明：只要 $a\in H^i(M)$、$i=0,1$，就对所有
$A\in H_2(M)\setminus\{0\}$ 和 $b,c\in H^*(M)$ 有
\[
  \operatorname{GW}_{A,3}^M(a,b,c)=0.
\]
若 $a\in H^0(M)$，这是命题~\ref{prop:7-5-6} 中（基本类）公理的特例。下面给出一个直接
论证。若 $\deg(a)\le1$，并且 $a,b,c$ 满足维数条件
\[
  \deg(a)+\deg(b)+\deg(c)=\dim M+2c_1(A),
\]
则
\[
  \deg(b)+\deg(c)>\dim M+2c_1(A)-2
  =\dim\mathcal M_{0,2}^*(A;J).
\]
因此，对 $J\in\mathcal J_{\mathrm{reg}}(M,\omega)$，\BookRecordIndex{T0283}求值映射
$\operatorname{ev}:\mathcal M_{0,2}^*(A;J)\to M\times M$ 避开表示 $b,c$ 的 Poincaré
对偶的两个循环的通有乘积 $X\times Y$，故不变量为零。

为证明 (iii)，注意当 $a,b,c\in H^*(M)$ 次数纯齐时，
$\operatorname{GW}_{A,3}^M(a,b,c)$ 在三个输入的置换下分次对称，即
\[
  \operatorname{GW}_{A,3}^M(a,b,c)
  =(-1)^{|a|(|b|+|c|)}\operatorname{GW}_{A,3}^M(b,c,a).
\]
由~\eqref{eq:11-1-9}，这意味着
\[
  \langle(a*b)_A,c\rangle
  =(-1)^{|a|(|b|+|c|)}\langle(b*c)_A,a\rangle
  =\langle a,(b*c)_A\rangle.
\]
这对所有 $A\in K^{\mathrm{eff}}$ 以及所有满足
$|a|+|b|+|c|=\dim M+2c_1(A)$ 的 $a,b,c\in H^*(M)$ 成立。因此 (iii) 由
\eqref{eq:11-1-6} 与~\eqref{eq:11-1-9} \BookRecordIndex{T0326}推出。命题~\ref{prop:11-1-11} 得证。
\end{proof}

命题~\ref{prop:11-1-11} 表明，小\BookRecordIndex{T0229}量子上同调环 $\operatorname{QH}^*(M;\Lambda)$ 是一个
\BookRecordIndex{T0019}Frobenius 代数。这类代数有许多引人注目的性质，并产生平坦联络和可积哈密顿系统。我们将在
第 11.5 节详细讨论。

\paragraph{例与习题。}

\begin{example}[复射影空间]\label{ex:11-1-12}
考虑\BookRecordIndex{T0358}辛流形 $M:=\mathbb{CP}^n$，赋予其标准复结构以及与 Fubini--Study 度量相联系的
凯勒 形式 $\omega$。令 $L\in H_2(\mathbb{CP}^n)$ 为由直线 $\mathbb{CP}^1$ 表示的
标准生成元。$\mathbb{CP}^n$ 的第一陈类满足
\[
  c_1(L)=n+1.
\]
因此 $\mathbb{CP}^n$ 单调，可以取 $\Lambda$ 为多项式环 $\mathbb Z[q]$，其中
\[
  \deg(q):=2(n+1).
\]
上同调 $H^*(\mathbb{CP}^n)$ 是由满足 $c(L)=1$ 的类
$c\in H^2(\mathbb{CP}^n)$ 生成的截断多项式环。因此，若 $k$ 为奇数，\BookRecordIndex{T0229}量子上同调群
$\operatorname{QH}^k(M)=\operatorname{QH}^k(M;\Lambda)$ 为零；对每个偶数 $k$，它有唯一
生成元 $c^iq^\ell$，其中
\[
  k=2\bigl(i+\ell(n+1)\bigr),\qquad0\le i\le n.
\]
这描述了 $\operatorname{QH}^*(M)$ 的加法结构。

为确定乘法结构，只需计算 $0\le i,j\le n$ 时的 $c^i*c^j$。由维数原因，除非 $d=0,1$，
不变量 $\operatorname{GW}_{dL,3}^{\mathbb{CP}^n}(a,b,c)$ 都为零。由命题
\ref{prop:11-1-11}(i)，$d=0$ 对应常值曲线并给出通常杯积。此外，
\[
  \operatorname{GW}_{L,3}^{\mathbb{CP}^n}(c^i,c^j,c^k)
  =
  \begin{cases}
  1,&i+j+k=2n+1,\\
  0,&\text{其他情形}.
  \end{cases}
\]
为看出这一点，注意与两个通有子空间相交的直线空间充满一个子空间，其中两个通有子空间
的复维数分别为 $n-i$ 和 $n-j$。这个子空间的维数为 $n-i+n-j+1$，且其中每一点恰有一条
这样的直线通过。因此，若 $c^k$ 的 Poincaré 对偶具有互补维数，恰有一条直线对
\BookRecordIndex{T0143}格罗莫夫--威滕不变量作出贡献。由此
\[
  c^i*c^j=
  \begin{cases}
  c^{i+j},&i+j\le n,\\
  c^{i+j-n-1}q,&n+1\le i+j\le2n.
  \end{cases}
\]
因此，对 $i\le n$，$c^i$ 等于 $c$ 的 $i$ 重量子积
$c*\cdots*c=:c^{*i}$，而量子环中的关系由唯一的新关系
\[
  c*c^n=q
\]
生成。这可以写成
\[
  \operatorname{QH}^*(\mathbb{CP}^n)
  =\frac{\mathbb Z[p,q]}{\langle p^{n+1}=q\rangle},
\]
其中 $p$ 对应于类 $c$。令 $q=0$，就恢复 $\mathbb{CP}^n$ 的通常上同调环。
\end{example}

下一个例子虽然简单，却很好地展示了系数环 $\Lambda$ 的选择对\BookRecordIndex{T0229}量子上同调环的代数结构有何
影响。尽管\BookRecordIndex{T0228}量子杯积由\BookRecordIndex{T0143}格罗莫夫--威滕不变量决定，而后者在 $\omega$ 的形变下不变（见
注~\ref{rem:7-1-11}），$\Lambda$ 的可用选择却依赖于\BookRecordIndex{T0363}辛形式的上同调类 $[\omega]$。
因而，有时可以用\BookRecordIndex{T0229}量子上同调获取关于 $[\omega]$ 的信息；见下面例 11.4.9。

\begin{example}[$S^2\times S^2$ 的乘积]\label{ex:11-1-13}
考虑流形 $M:=S^2\times S^2$，赋予标准辛结构
$\omega=\pi_1^*\sigma+\pi_2^*\sigma$。定义
\[
  A:=[S^2\times\mathrm{pt}],\qquad B:=[\mathrm{pt}\times S^2]
\]
为 $H_2(M)$ 中的类，并令 $a,b$ 是 $H^2(M)$ 的对偶基；于是
$a=\operatorname{PD}(B)$、$b=\operatorname{PD}(A)$。因此 $H^*(M)$ 有加法基
$1,a,b,ab:=a\smile b$。$M$ 中非平凡的三点不变量为
\[
  \operatorname{GW}_{A,3}^M(ab,a,a)=1,\qquad
  \operatorname{GW}_{B,3}^M(ab,b,b)=1,\qquad
  \operatorname{GW}_{A+B,3}^M(ab,ab,ab)=1.
\]
相应地，$a*b=a\smile b$（所以用 $ab$ 表示这个元素不会有歧义），其余量子积为
\[
\begin{gathered}
  a*a=1\otimes e^A,\qquad b*b=1\otimes e^B,\\
  a*ab=b\otimes e^A,\qquad b*ab=a\otimes e^B,\\
  ab*ab=1\otimes e^{A+B}.
\end{gathered}
\]
$a*a,a*b,b*b$ 的表达由结合律决定其他乘积（如 $a*ab$）；读者可以验证所有公式彼此
\BookRecordIndex{T0348}相容。

现在更仔细地考察系数环 $\Lambda$。\BookRecordIndex{T0390}有效锥为
\[
  K^{\mathrm{eff}}=\{d_1A+d_2B\mid d_1,d_2\ge0\},
\]
故 $\Gamma=\Gamma(M,\omega)$ 是 $K^{\mathrm{eff}}$ 的半群环，可等同于双变量多项式环。
下面分别考察 $\Lambda$ 为双变量多项式环、单变量多项式环和单变量 Laurent 级数域这三种
情形；所得\BookRecordIndex{T0229}量子上同调环的代数结构颇不相同。

先取 $\Lambda:=\mathbb C[q_1,q_2]$，其中同构 $\phi:\Gamma\to\Lambda$ 由
\[
  e^A:=q_1,\qquad e^B:=q_2
\]
给出（见例~\ref{ex:11-1-4}(ii)）。\sourcefootnote{原书此处回指例 11.1.4(iii)；
但此处的双变量多项式环正是例 11.1.4(ii) 中群环在选定整数基后的表示，
而 (iii) 定义的是 Novikov 完备化，故改正例号。}于是 $a*a=q_1$、$b*b=q_2$，从而
\[
  H^*(M;\mathbb C)\cong
  \frac{\mathbb C[p_1,p_2]}{\langle p_1^2=p_2^2=0\rangle},
  \qquad
  \operatorname{QH}^*(M;\Lambda)\cong
  \frac{\mathbb C[p_1,p_2,q_1,q_2]}
       {\langle p_1^2=q_1,\ p_2^2=q_2\rangle},
\]
其中
\[
  \deg(p_1)=\deg(p_2)=2,\qquad\deg(q_1)=\deg(q_2)=4.
\]
因此 $\operatorname{QH}^*(M;\Lambda)$ 同构于多项式环 $\mathbb C[p_1,p_2]$。后一种描述
隐藏了 $\Lambda$-模结构，却清楚显示这个环没有零因子。

由于处于单调情形，也可以取 $\Lambda:=\mathbb C[q]$ 为次数四的单变量多项式环，并令
$e^A:=e^B:=q$（例~\ref{ex:11-1-4}(i)）。这相当于在 $\mathbb C[q_1,q_2]$ 中把
$q_1,q_2$ 等同。此时 $a*a=b*b=q$，从而
\[
  \operatorname{QH}^*(M;\Lambda)\cong
  \frac{\mathbb C[p_1,p_2,q]}
       {\langle p_1^2=p_2^2=q\rangle}.
\]
这个环确有零因子，因为
$(p_1-p_2)(p_1+p_2)=p_1^2-p_2^2=0$。

最后令 $\Lambda$ 为复系数的 $q$-Laurent 级数域。单项式 $q$ 在 $\Lambda$ 中可逆，故
\[
  e_+:=\frac{q+ab}{2q},\qquad e_-:=\frac{q-ab}{2q}
\]
是\BookRecordIndex{T0229}量子上同调环中良定义的元素。关系式表明它们是两个相互交换的投影：
\[
  (e_\pm)^2=e_\pm,\qquad e_+e_-=e_-e_+=0.
\]
因此 $\operatorname{QH}^*(M;\Lambda)$ 是交换代数，并作为 $\Lambda$-向量空间由
$e_+,e_-,ae_+,ae_-$ 生成。由 $a^2=q$，可知它作为 $\Lambda$-代数同构于直和
\[
  e_+Q\oplus e_-Q,
\]
其中域 $Q$ 是给 $\Lambda$ 添加 $q$ 的平方根所得的扩张。这是 Abrams \srcite{7} 关于
某些单调流形以域为系数的\BookRecordIndex{T0229}量子上同调环结构之结果的一个特例。Entov--Polterovich
\srcite{105} 在哈密顿辛同胚群上构造量子同态时使用了这些观察。
\end{example}

\begin{exercise}\label{ex:11-1-14}
令
\[
  M:=\mathbb{CP}^2\#\overline{\mathbb{CP}}{}^{\,2}
\]
为 $\mathbb{CP}^2$ 的单点吹起。直接计算证明：存在系数环 $\Lambda$ 使
$\operatorname{QH}^*(M;\Lambda)$ 为域，当且仅当 $M$ 是\emph{大吹起}，即
\[
  3\omega(E)^2\ge\omega(L)^2,
\]
其中 $E$ 是例外除子的类，$L$ 是直线的类；比较 Ostrover \srcite{314, \S3.3}。现在，这一
现象已可用 $M$ 的标准环面结构中拉格朗日纤维的可移性性质完全解释：大吹起只有一条不可移
纤维，其他吹起则有两条。见 Fukaya--Oh--Ohta--Ono
\srcite{129, Theorem~1.12} 和 McDuff \srcite{275}。其他一些几何推论见 McDuff
\srcite{274} 的注 1.10 与 1.11。
\end{exercise}

\begin{example}[丛的\BookRecordIndex{T0246}模空间]\label{ex:11-1-15}
设 $\Sigma$ 是亏格 $g\ge2$ 的黎曼曲面，考虑其上具有固定次数一\BookRecordIndex{T0162}行列式线丛的稳定秩二全纯
向量丛之\BookRecordIndex{T0246}模空间 $M_\Sigma$。这是一个单连通的单调\BookRecordIndex{T0358}辛流形，且 $\pi_2=\mathbb Z$，因而其
\BookRecordIndex{T0229}量子上同调环良定义。Donaldson \srcite{83}、Muñoz \srcite{296} 和 Siebert--Tian
\srcite{377} 计算了这个环。详细描述这个例子会偏离主题太远，我们只指出其中几个有趣
特征。首先，Donaldson 的 $\mu$-映射给出同构
\[
  H_1(\Sigma;\mathbb Z)\longrightarrow H^3(M_\Sigma;\mathbb Z).
\]
因此奇数次上同调非常丰富，而且量子积在其中某些类上确实发生形变。换言之，存在
$a_1,a_2\in H^{\mathrm{odd}}(M_\Sigma)$、$a_3\in H^{\mathrm{ev}}(M_\Sigma)$ 以及
$0\ne A\in H_2(M_\Sigma)$，使得
\[
  \operatorname{GW}_{A,3}^{M_\Sigma}(a_1,a_2,a_3)\ne0.
\]
原则上没有理由要求奇数次类的\BookRecordIndex{T0143}格罗莫夫--威滕不变量消失，但要找出不消失的显式例子并不
容易；这是本书给出的唯一一个。精确公式例如见 Muñoz \srcite{296, \S4}。该文还很好地
展示了给出\BookRecordIndex{T0229}量子上同调环的表示与描述所有三点\BookRecordIndex{T0143}格罗莫夫--威滕不变量之间的区别；比较下面的
注 11.3.20。最后，该文对单调流形使用了稍有不同的\BookRecordIndex{T0229}量子上同调版本：在多项式系数环
$\mathbb Z[q]$ 中令 $q=1$。所得\BookRecordIndex{T0229}量子上同调环 $\operatorname{QH}^*(M)$ 与通常上同调环
$H^*(M;\mathbb Z)$ 含有完全相同的元素，但乘法发生形变，并且只模 $2N$ 分次，其中 $N$
是 $M$ 的\BookRecordIndex{T0418}最小陈数。在例~\ref{ex:8-6-8} 对 \BookRecordIndex{T0047}Seidel 表示的讨论中，我们实际上使用的就是
这个\BookRecordIndex{T0229}量子上同调版本。
\end{example}

\begin{example}\label{ex:11-1-16}
Donaldson 在 \srcite{83} 中计算了闭黎曼曲面上平坦 $\operatorname{SO}(3)$-丛之\BookRecordIndex{T0246}模空间
$M_\Sigma$ 的小\BookRecordIndex{T0229}量子上同调，所取的是最简单的非平凡情形：亏格为 $2$。由
Narasimhan--Seshadri 定理，该\BookRecordIndex{T0246}模空间同构于例~\ref{ex:11-1-15} 所讨论的\BookRecordIndex{T0246}模空间。它的实
维数为 $6$，\BookRecordIndex{T0418}最小陈数 $N=2$。计算的起点是如下观察：$M_\Sigma$ 辛同胚于
$\mathbb{CP}^5$ 中两个二次超曲面的横截交
\[
  M=X_1\cap X_2.
\]
用 $h\in H^2(M;\mathbb Z)$ 表示 $H^2(\mathbb{CP}^5;\mathbb Z)$ 的正生成元之拉回。则
$M$ 的陈类为
\[
  c_1(TM)=2h,\qquad c_2(TM)=3h^2,\qquad c_3(TM)=0.
\]
因此 $M$ 的 Euler 示性数为零；由 Lefschetz 超平面定理，其 Betti 数为
\[
  b_0=b_2=b_4=b_6=1,\qquad b_1=b_5=0,\qquad b_3=4.
\]
用 $h_0=1,h_2=h,h_4,h_6$ 表示偶数次生成元。于是 $h$ 与一个超平面截面对偶，$h_4$ 与
$M$ 中的一条直线对偶，$h_6$ 与一点对偶。因此
\[
  h\smile h=4h_4,\qquad h\smile h_4=h_6.
\]
（第二个等式显然，第一个来自恒等式 $h^3=4h_6$。）用
$(\alpha,\beta)\mapsto\langle\alpha,\beta\rangle$ 表示 $H^3(M)$ 上的相交配对，使得
\[
  \alpha\smile\beta=\langle\alpha,\beta\rangle h_6.
\]
Donaldson 证明，$M$ 中有 $4$ 条直线经过一个通有点，有 $2$ 条直线经过一对通有直线。由
这一几何信息出发，可用\BookRecordIndex{T0125}分裂公理计算类 $L,2L,3L$ 的\BookRecordIndex{T0143}格罗莫夫--威滕不变量，其中
$L\in H_2(M;\mathbb Z)$ 是正生成元，由 $M$ 中的直线表示。所得公式为
\[
\begin{gathered}
  \operatorname{GW}_{L}(h,h,h_6)=4,\qquad
  \operatorname{GW}_{L}(h,h_4,h_4)=2,\qquad
  \operatorname{GW}_{L}(\alpha,\beta,h_4)=-\langle\alpha,\beta\rangle,\\
  \operatorname{GW}_{2L}(h,h_6,h_6)=4,\qquad
  \operatorname{GW}_{2L}(h_4,h_4,h_6)=3,\qquad
  \operatorname{GW}_{3L}(h_6,h_6,h_6)=4.
\end{gathered}
\]
因此量子积为
\[
\begin{gathered}
  h*h=4h_4+4q,\qquad h*h_4=h_6+2hq,\qquad
  h*h_6=4h_4q+4q^2,\\
  h_4*h_4=2h_4q+3q^2,\qquad h_4*h_6=3hq^2,\qquad
  h_6*h_6=4h_4q^2+4q^3,
\end{gathered}
\]
并且对 $\alpha,\beta\in H^3(M)$，
\[
  h*\alpha=0,\qquad h_4*\alpha=-\alpha q,\qquad h_6*\alpha=0,\qquad
  \alpha*\beta=\langle\alpha,\beta\rangle(h_6-hq),
\]
其中 $\deg(q)=4$。所以 $M$ 的\BookRecordIndex{T0229}量子上同调环由一个次数二元素 $h$、次数三的类
$\alpha\in H^3(M)$ 以及一个次数四元素 $q$ 生成，并满足关系
\[
  h^4=16h^2q,\qquad h\alpha=0,\qquad
  \alpha\beta=\langle\alpha,\beta\rangle
  \left(\frac14h^3-4hq\right).
\]
而类 $h_4,h_6$ 可由 $h,q$ 表示为
\[
  h_4:=\frac14h^2-q,\qquad h_6:=\frac14h^3-3hq.
\]
这些公式与\BookRecordIndex{T0231}量子系数环的具体选择（例如取 $q$ 中的多项式、Laurent 多项式或 Laurent
级数）无关。
\end{example}

第 11.3 节还会给出更多\BookRecordIndex{T0229}量子上同调的例子。

\begin{exercise}\label{ex:11-1-17}
设 $a\in H^i(M)$、$b\in H^j(M)$，并写成
\[
  a*b=:\sum_A(a*b)_A\otimes e^A,
  \qquad (a*b)_A\in H^{i+j-2c_1(A)}(M).
\]
证明：$(a*b)_A$ 只有在 $A=0$ 且 $i+j=2n$ 时才可能是 $H^{2n}(M)$ 的非零元素。由此\BookRecordIndex{T0326}推出
\[
  \langle a*b,1\rangle=\langle a\smile b,1\rangle
  =\int_Ma\smile b.
\]
这给出了命题~\ref{prop:11-1-11}(iii) 一个特例的另一种证明。
\end{exercise}

现在推广到迭代乘积。令 $I:=\{1,\ldots,k\}$，考虑第~\ref{sec:7-3} 节定义的带\BookRecordIndex{T0150}固定标记点
的\BookRecordIndex{T0143}格罗莫夫--威滕不变量
\[
  \operatorname{GW}_{A,k}^{M,I}(a_1,\ldots,a_k),
  \qquad a_i\in H^*(M).
\]
（分裂）公理蕴含，这些不变量由三点不变量决定。下一个习题精确说明这一点，并证明它们
对应于迭代量子积。Entov 在 \srcite{104} 中以特殊酉矩阵乘积之特征值所满足的 ABW
不等式为语言，对格拉斯曼流形的这些数给出了一个优美解释。

\begin{exercise}[迭代量子积]\label{ex:11-1-18}
设 $(M,\omega)$ 是闭\BookRecordIndex{T0058}半正\BookRecordIndex{T0358}辛流形，$\Lambda=\Lambda_\omega$ 是
例~\ref{ex:11-1-4}(iii) 的整数系数 Novikov 环，即 $R=\mathbb Z$。
\sourcefootnote{原书把 $\Lambda_\omega$ 回指为例 11.1.4(iv)；但
$\Lambda_\omega$ 明确定义于例 11.1.4(iii)，而 (iv) 定义的是\BookRecordIndex{T0270}普适 Novikov 环，故改正
例号。}
回忆注~\ref{rem:11-1-7}：此时每个\BookRecordIndex{T0229}量子上同调类
$a\in\operatorname{QH}^*(M;\Lambda)$ 都形如
\[
  a=\sum_Aa_A\otimes e^A,
\]
其中求和遍历所有同调类 $A\in H_2(M)$，并且对每个 $c\in\mathbb R$ 都有
\[
  \#\{A\in H_2(M)\mid a_A\ne0,\ \omega(A)\le c\}<\infty.
\]
在这一情形，自然可以把带\BookRecordIndex{T0150}固定标记点的\BookRecordIndex{T0143}格罗莫夫--威滕不变量之定义延拓如下：令
$I_k:=\{1,\ldots,k\}$，得到同态
\[
  \operatorname{GW}_{A,k}^{M,I_k}:
  \underbrace{\operatorname{QH}^*(M;\Lambda)\times\cdots\times
  \operatorname{QH}^*(M;\Lambda)}_{k\ \text{个}}
  \longrightarrow\mathbb Z,
\]
其公式为
\[
  \operatorname{GW}_{A,k}^{M,I_k}(a_1,\ldots,a_k)
  :=\sum_{A_1,\ldots,A_k}
  \operatorname{GW}_{A-A_1-\cdots-A_k,k}^{M,I_k}
  (a_{1,A_1},\ldots,a_{k,A_k}),
\]
其中
\[
  a_i=\sum_Ba_{i,B}\otimes e^B\in\operatorname{QH}^*(M;\Lambda).
\]
依定义，$\operatorname{GW}_{A,k}^{M,I_k}$ 下降到
\[
  \operatorname{QH}^*(M;\Lambda)\otimes_\Lambda\cdots
  \otimes_\Lambda\operatorname{QH}^*(M;\Lambda).
\]
证明
\[
  \operatorname{GW}_{A,k}^{M,I_k}(a_1,\ldots,a_k)
  =\operatorname{GW}_{A,1}^{M,I_1}(a_1*\cdots*a_k)
  =\int_M(a_1*\cdots*a_k)_A,
\]
其中 $\operatorname{GW}_{A,1}^{M,I_1}$ 是习题~\ref{ex:7-3-4} 的一点不变量，按上述方式
延拓到 $\operatorname{QH}^*(M;\Lambda)$。因此迭代积可写成
\[
  a_1*\cdots*a_k=\sum_A(a_1*\cdots*a_k)_A\otimes e^A,
\]
其中上同调类 $(a_1*\cdots*a_k)_A\in H^*(M)$ 由
\[
  \langle(a_1*\cdots*a_k)_A,c\rangle
  :=\operatorname{GW}_{A,k+1}^{M,I_{k+1}}(a_1,\ldots,a_k,c),
  \qquad c\in H^*(M),
\]
定义。\emph{提示：}用归纳法，并在所有\BookRecordIndex{T0066}标记点都固定的情形使用定理~\ref{thm:7-5-10}
的\BookRecordIndex{T0254}粘合公式。
\end{exercise}

\begin{exercise}[Künneth 公式]\label{ex:11-1-19}
假定 $(M_1,\omega_1)$ 与 $(M_2,\omega_2)$ 以相同常数单调，因而
$(M_1\times M_2,\pi_1^*\omega_1+\pi_2^*\omega_2)$ 也单调。令
$A_i\in H_2(M_i)$，$a_i,b_i,c_i\in H^*(M_i)$（$i=1,2$）。

以下定义\sourcefootnote{原书把 $a,b,c$ 中两因子拉回类之间的杯积都误印成加号；但紧接着的
\BookRecordIndex{T0143}格罗莫夫--威滕不变量乘积分解、多线性和次数条件都要求这里取\BookRecordIndex{T0331}外杯积，故改正。原书随后
的四点公式又直接使用了未定义的 $d_1,d_2,d$；按前三组变量的同一模式补出其定义。}
采用\BookRecordIndex{T0331}外杯积，并记
\[
\begin{gathered}
  A:=\iota_{1*}A_1+\iota_{2*}A_2,\qquad
  a:=\pi_1^*a_1\smile\pi_2^*a_2,\\
  b:=\pi_1^*b_1\smile\pi_2^*b_2,\qquad
  c:=\pi_1^*c_1\smile\pi_2^*c_2.
\end{gathered}
\]
证明
\[
  \operatorname{GW}_{A,3}^{M_1\times M_2}(a,b,c)
  =\operatorname{GW}_{A_1,3}^{M_1}(a_1,b_1,c_1)\,
   \operatorname{GW}_{A_2,3}^{M_2}(a_2,b_2,c_2).
\]
由此\BookRecordIndex{T0326}推出：若像例~\ref{ex:11-1-4}(i) 那样取单变量多项式环 $\mathbb Z[q]$ 为系数，则
\[
  \operatorname{QH}^*(M_1\times M_2;\mathbb Z[q])
  \cong
  \operatorname{QH}^*(M_1;\mathbb Z[q])
  \otimes_{\mathbb Z[q]}
  \operatorname{QH}^*(M_2;\mathbb Z[q]).
\]
所以小\BookRecordIndex{T0229}量子上同调的乘积公式非常简单。\BookRecordIndex{T0083}大量子上同调的情形有趣得多，因为其中涉及
$k>3$ 的 $k$ 点不变量。再取 $d_i\in H^*(M_i)$，并令
$d:=\pi_1^*d_1\smile\pi_2^*d_2$。
证明在上述记号下，当 $k=4$ 时有恒等式
\[
\begin{aligned}
  \operatorname{GW}_{A,4}^{M_1\times M_2}(a,b,c,d)
  &=
  \operatorname{GW}_{A_1,4}^{M_1}(a_1,b_1,c_1,d_1)\,
  \operatorname{GW}_{A_2,4}^{M_2,\{1,2,3,4\}}(a_2,b_2,c_2,d_2)\\
  &\quad+
  \operatorname{GW}_{A_1,4}^{M_1,\{1,2,3,4\}}(a_1,b_1,c_1,d_1)\,
  \operatorname{GW}_{A_2,4}^{M_2}(a_2,b_2,c_2,d_2).
\end{aligned}
\]
一般 Künneth 公式的陈述见 Kaufmann \srcite{207}。
\end{exercise}

\begin{remark}[\BookRecordIndex{T0230}量子同调]\label{rem:11-1-20}
本节经由\BookRecordIndex{T0145}格罗莫夫--威滕伪循环的 Poincaré 对偶定义了\BookRecordIndex{T0228}量子杯积。采用上同调往往很合适，
但有时从同调角度思考更有用，因为它与几何的联系更直接。$M$ 的同调上的量子积是相交积的
形变，定义为\BookRecordIndex{T0228}量子杯积的 Poincaré 对偶。因此，两个次数纯齐的同调类
$\alpha,\beta\in H_*(M)$ 的量子积为
\begin{equation}\label{eq:11-1-11}
  \alpha*\beta:=
  \sum_A\sum_{\nu,\mu}
  \operatorname{GW}_{A,3}^M(\alpha,\beta,\varepsilon_\nu)
  g^{\nu\mu}\varepsilon_\mu\otimes e^{-A},
\end{equation}
其中 $\varepsilon_\nu$ 构成 $H_*(M)$ 的一个基，$g^{\nu\mu}$ 是相交矩阵
$g_{\nu\mu}:=\varepsilon_\nu\cdot\varepsilon_\mu$ 的逆矩阵。约定
$\deg(e^{-A}):=-2c_1(A)$ 后，量子积的次数满足
\[
  \deg(\alpha*\beta)=\deg(\alpha)+\deg(\beta)-2n,
\]
而首项（即 $A=0$ 时 $e^{-A}$ 的系数）就是相交积 $\alpha\cdot\beta$。还可以写成
$\alpha*\beta=\sum_A(\alpha*\beta)_A\otimes e^{-A}$，其中同调类
$(\alpha*\beta)_A$ 是 $\alpha$ 与 $\beta$ 的一种“增厚”相交；正如
注~\ref{rem:11-1-9} 更精确地说明的，它由所有同时与 $\alpha,\beta$ 相交的 $A$-曲线之并
表示。

为了使这个定义有意义，需要解释可能为无穷和的表达式
\begin{equation}\label{eq:11-1-12}
  \alpha=\sum_A\alpha_A\otimes e^{-A},
\end{equation}
其中 $\alpha_A\in H_{k-2c_1(A)}(M)$ 对每个 $c\in\mathbb R$ 满足
\[
  \#\{A\mid\alpha_A\ne0,\ \omega(-A)\ge c\}<\infty.
\]
\BookRecordIndex{T0230}量子同调应具有如下性质。第一，应存在 Poincaré 对偶同构
$\operatorname{PD}:\operatorname{QH}^k\to\operatorname{QH}_{2n-k}$，其在
形如~\eqref{eq:11-1-7} 的类上限制为
\[
  a=\sum_Aa_A\otimes e^A
  \longmapsto
  \operatorname{PD}(a):=\sum_A\operatorname{PD}(a_A)\otimes e^{-A}.
\]
第二，应存在一个（最好非退化的）配对
$\operatorname{QH}^*\times\operatorname{QH}_*\to R$，对分别具有
\eqref{eq:11-1-7} 与~\eqref{eq:11-1-12} 形式的 $a,\alpha$ 满足
\[
  \langle a,\alpha\rangle=\sum_A\int_{\alpha_{-A}}a_A.
\]

为严格化上述说法，假定 $(\Lambda,\phi,\iota)$ 是定义~\ref{def:11-1-3} 中以底环 $R$
为系数的\BookRecordIndex{T0231}量子系数环。定义\emph{\BookRecordIndex{T0108}对偶环} $\check\Lambda$：它作为环等于 $\Lambda$，但分次
反向，即
\[
  \check\Lambda^k:=\Lambda^{-k}.
\]
令 $r:\Lambda\to\check\Lambda$ 表示这个显然的（反转次数的）同构，并定义
\[
  \check\phi:=r\circ\phi:\Gamma(M,\omega)\to\check\Lambda,\qquad
  \check\iota:=\iota\circ r^{-1}:\check\Lambda\to R.
\]
三元组 $(\check\Lambda,\check\phi,\check\iota)$ 称为\emph{\BookRecordIndex{T0109}对偶量子系数环}。使用记号
\[
  e^{-A}:=r(e^A)=\check\phi(\delta_A)\in\check\Lambda
\]
很方便。现在定义
\[
  \operatorname{QH}_*(M;\check\Lambda)
  :=H_*(M)\otimes_{\mathbb Z}\check\Lambda.
\]
与注~\ref{rem:11-1-7} 一样，只要函数
$K^{\mathrm{eff}}\to H_*(M):A\mapsto\alpha_A$ 满足
$\sup_{\alpha_A\ne0}|c_1(A)|<\infty$，和~\eqref{eq:11-1-12} 就是
$\operatorname{QH}_*(M;\check\Lambda)$ 中良定义的元素。因此，公式~\eqref{eq:11-1-11}
在 $\operatorname{QH}_*(M;\check\Lambda)$ 上定义量子相交积。

此外，有 Poincaré 对偶同构
\[
\begin{aligned}
  \operatorname{QH}^k(M;\Lambda)&\longrightarrow
  \operatorname{QH}_{2n-k}(M;\check\Lambda),\\
  a=\sum_i a_i\otimes\lambda_i&\longmapsto
  \operatorname{PD}(a):=\sum_i\operatorname{PD}(a_i)\otimes r(\lambda_i),
\end{aligned}
\]
以及配对
\[
  \operatorname{QH}^k(M;\Lambda)\times
  \operatorname{QH}_k(M;\check\Lambda)\longrightarrow R,
  \qquad(a,\alpha)\longmapsto\langle a,\alpha\rangle,
\]
其中
\begin{equation}\label{eq:11-1-13}
  \langle a,\alpha\rangle
  :=\sum_{i,j}\langle\lambda_i,\check\lambda_j\rangle
  \int_{\alpha_j}a_i
\end{equation}
对 $a=\sum_i a_i\otimes\lambda_i$ 和
$\alpha=\sum_j\alpha_j\otimes\check\lambda_j$ 定义，而
\[
  \langle\lambda,\check\lambda\rangle
  :=\iota\bigl(\lambda r^{-1}(\check\lambda)\bigr).
\]
读者可以验证这些结构满足上述要求。在 Poincaré 对偶下，\BookRecordIndex{T0229}量子上同调与\BookRecordIndex{T0230}量子同调之间的配对
同构于 $\operatorname{QH}^*(M;\Lambda)$ 上的双线性形式~\eqref{eq:11-1-6}。它未必
非退化；当配对
\[
  \Lambda\times\check\Lambda\longrightarrow R,\qquad
  (\lambda,\check\lambda)\longmapsto
  \langle\lambda,\check\lambda\rangle
  :=\iota\bigl(\lambda r^{-1}(\check\lambda)\bigr)
\]
非退化时，它才非退化。
\end{remark}

\begin{exercise}\label{ex:11-1-21}
证明：对 Laurent 多项式环 $\Lambda=R[q,q^{-1}]$ 以及 $q$ 中的 Laurent 级数环，配对
$\Lambda\times\check\Lambda\to R$ 非退化。再证明：对多项式环 $\Lambda=R[q]$ 以及
$q$ 中的形式幂级数环，它是退化的。
\end{exercise}

\begin{exercise}\label{ex:11-1-22}
证明配对~\eqref{eq:11-1-13} 满足恒等式
\[
  \langle a,\operatorname{PD}(b)\rangle
  =\langle a*b,1\rangle
  =\langle a,b\rangle,
  \qquad a,b\in\operatorname{QH}^*(M).
\]
这里 $\langle a,\operatorname{PD}(b)\rangle$ 由~\eqref{eq:11-1-13} 定义，
$\langle a,b\rangle$ 由~\eqref{eq:11-1-6} 定义；Poincaré 对偶同构中的符号要选成使
这个恒等式在通常（上）同调上成立。
\end{exercise}
